\section{Experiments}
\label{sec:exp}

The FLAN motivation studies in \figref{fig:acl-motivation} examine a fixed shared-prefix protocol with eight clients and ranks 16 and 64.
We report available federated LoRA baselines on NI and GLUE for context.
\TODO{Add the method comparison after the collaboration metric is defined and the corresponding experiments are complete.}

\subsection{Datasets}
\label{sec:exp:datasets}
We use the NI and GLUE Dirichlet $\alpha=0.5$ partitions, each with 16 clients.
The total rank budgets are 16, 32, and 64 for six, five, and five clients, respectively.
Each client has 200 test examples; the training and validation allocations for the reported methods are specified below.

\subsection{Evaluation}
\label{sec:exp:eval}
We report client-macro ROUGE-L (R-L) for Super-NaturalInstructions (NI) and GLUE.
For both datasets, each client generates answers on its test set, and the resulting client scores are averaged equally.
GLUE follows the project's causal generation protocol, so its reported metric is ROUGE-L rather than task-specific classification accuracy.

\subsection{Training Setup}
\label{sec:exp:environments}
The runs use Qwen3.5-0.8B and seed 615, with 10 main local optimizer steps per round and an effective batch size of eight.
We evaluate checkpoint R9 after ten communication rounds.
Method-specific optimizer phases and inference procedures incur additional computation beyond these main training steps.

\subsection{Main Results}
\label{sec:exp:results}
\tabref{tab:exp-ni-dir05} reports the available NI baselines.
\tabref{tab:exp-glue-dir05} reports the available GLUE baselines.

\begin{table}[t]
\centering
\small
\setlength{\tabcolsep}{4pt}
\caption{Client-macro ROUGE-L (\%) on Super-NaturalInstructions under Dirichlet $\alpha=0.5$.
The highest and second-highest scores are shown in bold and underlined, respectively.}
\label{tab:exp-ni-dir05}
\begin{tabular}{lr}
\toprule
\textbf{Method} & \textbf{ROUGE-L (\%)} \\
\midrule
LocalOnly                          & \textbf{57.5208} \\
FlexLoRA                           & \underline{54.7099} \\
FLoRA                              & 51.0599 \\
HetLoRA                            & 51.8617 \\
PF$^2$LoRA                          & 48.4953 \\
pFedLoRA                           & 51.3480 \\
FedRoRA                            & 53.6302 \\
Het-CPFLoRA                        & 54.0703 \\
\bottomrule
\end{tabular}
\end{table}

\begin{table}[t]
\centering
\small
\setlength{\tabcolsep}{4pt}
\caption{Client-macro baseline ROUGE-L (\%) on GLUE under Dirichlet $\alpha=0.5$.
The highest and second-highest scores are shown in bold and underlined, respectively.}
\label{tab:exp-glue-dir05}
\begin{tabular}{lr}
\toprule
\textbf{Method} & \textbf{ROUGE-L (\%)} \\
\midrule
LocalOnly                          & \textbf{79.9531} \\
FlexLoRA                           & \underline{79.9010} \\
FLoRA                              & 77.5885 \\
FedRoRA                            & 77.1719 \\
\bottomrule
\end{tabular}
\end{table}

\fakeparagraph{Comparison Scope}
PF$^2$LoRA uses shared rank 4 and private rank 12, 28, or 60 to match each client's total rank budget.
Het-CPFLoRA uses shared rank 12, 24, or 48 and its paper dynamic weighted inference protocol.
For NI, LocalOnly, FlexLoRA, FLoRA, and FedRoRA train on 936 examples per client with 64 held out, whereas HetLoRA, PF$^2$LoRA, pFedLoRA, and Het-CPFLoRA use all 1000 training examples.
The results therefore retain method-specific training-data and computation budgets.

\fakeparagraph{Observed Results}
Among the reported baselines, LocalOnly has the highest ROUGE-L on NI and GLUE.
